\chapter{Deployment and Reconfiguration Strategies on the Edge-Cloud Continuum}
\label{chap:deployment-and-reconfiguration}

Every correct deployment of a well-formed macro-program computes the same thing (\cref{thm:deployment-independence}).
%
The assignment of components to hosts is therefore free: it can be chosen, and later revised, on non-functional grounds alone, without touching the application logic (\cref{sec:pulverization-discussion}).
%
The language side is settled too:
aggregate computing, choreographic programming, and multitier programming can all produce the deployable components the model expects (\cref{chap:language-and-coordination}).
%
Neither result says who uses that freedom, or how a deployment is kept good while the infrastructure changes under it.

Any answer has to settle two decisions.
%
The first is what the deployment is for.
It fixes which quantity a candidate deployment should improve -- energy drawn, battery remaining, or carbon emitted -- and how much latency or extra communication that improvement may cost.
%
The second is when a running deployment is revisited:
continuously, in reaction to a local signal, or at fixed intervals, by recomputing the assignment against the infrastructure's current state.

This chapter evaluates three policies, each settling the two decisions differently.
%
The first is a single local rule:
a device offloads one resource-heavy component to a nearby host as soon as its own battery drops below a threshold, and moves it again if that host becomes unreachable (\cref{sec:battery-triggered-offloading}).
Nothing is searched, and nothing is stated beyond the threshold itself.
%
The second is self-organising as well, but the devices decide together.
Infrastructural devices compete for the components application devices would rather not run locally, resizing through a collectively built field the region each one leads,
so a device is offloaded to whichever infrastructural device combines proximity with spare capacity (\cref{sec:global-self-organising-reconfiguration}).
%
The third states the decision declaratively.
An objective over energy and carbon is given once, together with the hardware and latency requirements of every component,
and a planner searches the admissible deployments for one that meets it, repeating the search as the network drifts (\cref{sec:declarative-green-planning}).

The three are compared on the same terms at the end of the chapter (\cref{sec:deployment-reconfiguration-discussion}):
how much they improve the quantity they target -- energy, battery life, carbon, or the share of devices able to offload at all --
and what that improvement costs in extra hops, extra latency, or a slower transient before it appears.

\section{A Self-organising Offloading Rule Triggered by Battery Level}
\label{sec:battery-triggered-offloading}

The evaluation uses the rescue scenario of \cref{fig:rescue-case-study}, reduced to the two components the offloading decision affects:
Emergency Service and Steering Service, both collective~\cite{DBLP:conf/acsos/FarabegoliVC24}.
%
Biometrics Acquisition, Navigation System, and Data Analysis are left out, since none of them is a candidate for offloading.
%
Emergency Service is the expensive component: it locates an emergency from the readings of neighbouring devices.
Steering Service only guides a rescuer once a location is known, at a negligible computational cost.
%
The rule evaluated here therefore moves Emergency Service alone, and Steering Service always stays at its owner.

In the monolithic deployment both components run at their owner.
This is the baseline already used to check deployment independence (\cref{subsec:pulverization-validation}).
%
The deployment evaluated here offloads Emergency Service instead.
%
Each device starts a run with a battery level drawn at random.
As soon as that level falls below 30\%, the device forwards its Emergency Service instance to the nearest edge server acting as its surrogate (\cref{subsec:forwarding-chains}),
and switches to another surrogate if it loses the connection to the current one.
%
No other condition drives the decision: not load, not the state of the network elsewhere, only the device's own battery.
%
The rule is self-organising in the narrowest sense:
each device decides for itself, from one locally observed number, without negotiating and without knowing what the other devices do.

Reading the result requires a power model for the two device classes.
%
The reference values come from Shao and Brooks's per-instruction energy measurements~\cite{DBLP:conf/islped/ShaoB13}.
%
The smartphone is a Snapdragon 888, with a thermal design power of 5 W and an energy cost of 0.55 J per instruction.
The edge server is a Xeon Platinum P-8124, with a thermal design power of 220 W and an energy cost of 10 J per instruction.
%
A 4500 mAh battery at 3.7 V fixes the energy budget each device starts a run with.

The three network sizes used to validate deployment independence -- 50, 100, and 200 mobile devices -- are reused here (\cref{subsec:pulverization-validation}).
%
The simulations run again in Alchemist~\cite{DBLP:journals/jos/PianiniMV13}, with eight edge servers and one cloud instance available as surrogates,
under the mobility pattern already described for the functional validation.
%
\Cref{fig:power-consumption} compares the monolithic deployment against the one that offloads Emergency Service under the battery rule.

\begin{figure}[t]
    \centering
    \includegraphics[width=\textwidth]{power-consumption.pdf}
    \caption[Battery level and message overhead under a monolithic and an offloaded deployment]{Average battery level (left) and cumulative messages exchanged (right) by application devices, under a monolithic and an offloaded deployment. Data from~\cite{DBLP:conf/acsos/FarabegoliVC24}.}
    \label{fig:power-consumption}
\end{figure}

The offloaded deployment leaves devices with more battery at the end of the run.
%
It exchanges slightly more messages in return, as the right-hand chart of \cref{fig:power-consumption} shows:
reaching Emergency Service at its surrogate and returning its result takes hops that a local execution does not need.
%
The cost is not paid in energy, since the offloaded devices are the ones with more battery left.
It is paid in bandwidth, in the extra messages the forwarding chain to the surrogate requires.

The rule obtains this result from the cheapest possible statement: one threshold, one component, and no coordination between devices.
%
It does not generalise beyond the signal and the component it was written for.
%
A second bottleneck, or a second component worth offloading under different conditions, would need a second threshold chosen by hand.

\section{A Field-based Reconfiguration Policy over Self-organising Coordination Regions}
\label{sec:global-self-organising-reconfiguration}

This section answers the same question as \cref{sec:battery-triggered-offloading}:
which infrastructural device should run the component an application device would rather not run locally.
%
The answer is no longer read off a single, locally observed number.

Every infrastructural device seeds a potential field weighted by its own current load.
%
A lightly loaded device's field reaches further into the network than a heavily loaded one's,
and each application device is claimed by the field that reaches it with the lowest potential~\cite{DBLP:journals/iot/FarabegoliPCV24}.
%
The construction is the self-organising coordination regions pattern already met in \cref{subsec:pulverization-validation}~\cite{DBLP:journals/fgcs/PianiniCVN21}, with one difference.
%
Leaders are not elected at a fixed grain: every infrastructural device is already a leader,
and its region-forming field encodes its own current load rather than distance alone.
A leader carrying more work therefore commands a smaller region.

A converge-cast brings each region's candidate components to its leader.
%
The leader accepts as many as its own capacity allows and casts the decision back over the same region,
so every application device either keeps its component locally or runs it on the infrastructural device that claimed it.
%
The regions resize themselves as load shifts, so the decision follows the whole network's topology instead of one local signal.
%
The policy still runs over the same infrastructure/application-device split the middleware maintains (\cref{chap:pulverization}).

The policy is evaluated against a local baseline that pins every application device to its nearest infrastructural device by hop count, whatever that device's load.
%
The networks hold a thousand hosts each, arranged as a scale-free and as a lobster, tree-like topology.
%
Two scenarios are simulated:
a variable load imposed externally on the infrastructural devices,
and a degradation and recovery in which batches of infrastructural devices are switched off and then progressively reactivated.
%
Quality of service is the share of application devices that manage to offload their component at all.

Across both scenarios and both topologies, the field-based policy reaches a higher quality of service than the local baseline once its regions have stabilised.
%
Stabilisation takes time, and quality of service can dip below the baseline's during that transient before recovering past it.
%
The policy also degrades and recovers more gracefully than the baseline when infrastructural devices fail.
%
Both advantages are larger in the scale-free topology than in the lobster one:
removing a node from a lobster network is far more likely to segment it,
cutting off the alternative routes the field-based policy needs to redirect load elsewhere~\cite{DBLP:journals/iot/FarabegoliPCV24}.

\section{A Declarative Planner for Green Deployments}
\label{sec:declarative-green-planning}

Declarative continuous reasoning was introduced above as Forti and Brogi's incremental repair of a stated deployment goal (\cref{subsec:declarative-deployment})~\cite{DBLP:journals/logcom/FortiBB22}.
%
The planner presented here follows the same idea, worked out over the deployment units defined in \cref{chap:pulverization}~\cite{DBLP:conf/coordination/BrogiCFFV25}.

Every deployment unit declares its hardware requirements and the maximum latency it tolerates towards the unit of the same macro-program it depends on most closely.
%
Every candidate host declares its available and total capacity, the power-usage effectiveness of running work on it, and the mix of energy sources that supplies it, with the associated carbon intensity.
%
A Prolog-based generator searches this space for an assignment of every deployment unit to a host that meets both kinds of requirement.
The requirements are stated as facts rather than computed by a bespoke solver, and the search is ordinary logic-program resolution.

The assignment problem is the same combinatorial one that is already too costly to solve exactly at scale (\cref{subsec:constraint-based-placement}).
%
The generator therefore tries candidate hosts from the least to the most carbon-intensive, reaching a low-carbon assignment quickly instead of exploring the space exhaustively.
%
For every candidate it considers, the same declarative program estimates the resulting energy draw and carbon emissions
from each host's load before and after the placement, its power-usage effectiveness, and the carbon intensity of its energy mix.

The generator is invoked at a fixed interval from inside Alchemist~\cite{DBLP:journals/jos/PianiniMV13}, the same simulator used to validate the model (\cref{subsec:pulverization-validation}).
%
Each time it runs, it receives the infrastructure's current state and returns a new assignment,
so the running deployment is periodically replaced by the one the search currently favours.
%
This is the repeated declarative repair of \cref{subsec:declarative-deployment}, driven here by an explicit energy and carbon objective rather than by an incremental reasoner reacting to a stated goal.

The published planner states its requirements over five deployment units per device: sensing, actuation, behaviour, state, and communication.
%
This is the fixed partitioning of a logical device that precedes the present model (\cref{sec:pulverization-motivation}).
%
The latency bound is stated towards the state unit, since the other four were designed to depend on it.

The search itself is not tied to that count or to a single hub.
%
A hardware requirement and a latency bound towards whatever a unit depends on are the kind of per-component declaration the bindings of \cref{subsec:macroprogram-dag} already make explicit for a macro-program of any size,
and the same Prolog predicates would run unchanged over a \ac{DAG} of arbitrarily many components.
%
What is evaluated below is the five-unit instantiation, not a macro-program pulverised into such a \ac{DAG}.
%
The mechanism generalises; the reported numbers do not speak on their own for deployments finer than five units per device.

The evaluation is synthetic.
%
A variable number of physical devices, each able to run every component locally, are placed uniformly at random in a two-dimensional area whose size grows with the device count.
%
Devices connect to one another within a range that also grows with the device count, a stand-in for a short-range radio, and all of them connect to a single cloud instance.
%
Each device's mix of energy sources follows a jittered sine wave, capped at 90\% green energy for devices and 50\% for the cloud.
%
The sine wave stands in for the day-night cycle in a grid's carbon intensity, the time-varying supply behind Wiesner et al.'s carbon-aware scheduling (\cref{subsec:energy-aware-orchestration})~\cite{DBLP:conf/middleware/WiesnerBSGT21}.
%
Each run simulates 720 simulated minutes of device movement, repeated ten times per configuration to average out random placement and mobility.
%
Two deployments are compared over the same runs.
The peer-to-peer baseline keeps every unit at its owner for the whole run, which is the monolithic deployment in this thesis's terms.
The other is recomputed by the generator every 30 simulated minutes.

\begin{figure}[t]
    \centering
    \includegraphics[width=\textwidth]{green-planner-carbon-energy.pdf}
    \caption[Energy consumption and carbon footprint under a peer-to-peer baseline and the declarative planner]{Overall energy consumption (orange) and carbon footprint (blue) across network sizes, comparing the peer-to-peer baseline (top row) with the deployment recomputed every 30 simulated minutes by the declarative planner (bottom row). Shaded bands show one standard deviation across ten repetitions. Data from~\cite{DBLP:conf/coordination/BrogiCFFV25}.}
    \label{fig:green-planner-carbon-energy}
\end{figure}

At every network size tested, the planner's deployment draws roughly a third of the energy the peer-to-peer baseline draws on the same network:
under 2 kWh for 50 devices, rising to 3.5 kWh for 100.
%
Its carbon footprint follows the same pattern (\cref{fig:green-planner-carbon-energy}).
%
The baseline's carbon footprint tracks the sinusoid of the energy mix exactly, because its deployment never changes and cannot exploit a momentarily greener mix elsewhere in the network.
%
The planner's footprint departs from that curve.
%
Recomputing the assignment every 30 minutes steers deployment units towards the host whose energy mix is least carbon-intensive at the time.
This is, for a pulverised deployment, the time-varying objective described for a single scheduled job in \cref{subsec:energy-aware-orchestration}.

\begin{figure}[t]
    \centering
    \includegraphics[width=\textwidth]{green-planner-latencies.pdf}
    \caption[Intra-component and inter-device latency under a peer-to-peer baseline and the declarative planner]{Intra-component latency (blue), the average delay a device experiences reaching its own deployment units wherever they have been placed, and inter-device latency (orange), the delay between the communication units of neighbouring devices, across network sizes, for the peer-to-peer baseline and the declarative planner. Shaded bands show one standard deviation across ten repetitions. Data from~\cite{DBLP:conf/coordination/BrogiCFFV25}.}
    \label{fig:green-planner-latencies}
\end{figure}

The energy reduction of \cref{fig:green-planner-carbon-energy} is not free.
\Cref{fig:green-planner-latencies} shows where it is spent.

Intra-component latency is the delay a device experiences reaching its own deployment units, now spread across more than one host.
%
It rises well above the near-zero, in-memory cost of the peer-to-peer baseline, which is the expected price of leaving units off the device that owns them.
%
The increase stays within the same order of magnitude as the network grows,
so the search keeps a device's units close together rather than scattering them further as the problem gets larger.

Inter-device latency is the delay between the communication units of neighbouring devices, and it is the one that matters for the collective computation itself.
%
It rises only slightly, even though nothing in the objective rewards keeping communicating units nearby.
%
Concentrating a device's units on a few nearby hosts keeps neighbouring devices' communication units close to one another as well.
%
The search itself scales:
computing a new assignment takes about 100 ms whatever the network size, well inside the 30-minute interval between replanning rounds~\cite{DBLP:conf/coordination/BrogiCFFV25}.

\section{Discussion}
\label{sec:deployment-reconfiguration-discussion}

All three evaluations show the same thing:
revising who runs a component at runtime reduces the quantity the policy targets, at a bounded secondary cost.
%
The offloading rule pays in messages, the planner in latency, and the field-based policy in a stabilisation transient.
%
None of them trades one non-functional gain for another outright.

None of them removes the work of deciding what counts as a good deployment either.
%
The offloading rule needs a person to fix a battery threshold and to choose, in advance, the one component worth offloading.
%
The field-based policy needs a person to decide that load is the quantity its regions compete over.
%
The planner needs a person to state the energy and carbon objective, the constraints the search must respect, and how often the search is repeated.
%
Every one of these policies is written by hand.

They differ in where that judgement is exercised and how far it reaches.
%
A single threshold says nothing about a second bottleneck, a second component, or a trade-off between two competing objectives.
Extending the rule means writing a new rule.
%
The field-based policy reasons over the whole network's topology and recovers from a failed host better than the other two.
That recovery depends on the topology it runs on, a dependence the threshold rule and the planner never face, since each was tested over a single topology.
%
The planner's objective covers several concerns at once: energy, carbon, and a latency bound enforced as a hard constraint.
Someone still had to write that objective and those constraints beforehand, and a further concern has to be stated by hand as well.

The next chapter drops that requirement (\cref{chap:heterogeneous-gnn}).
%
Instead of a rule or an objective fixed in advance,
a policy is trained to weigh the same non-functional concerns from experience,
under a condition none of the three policies of this chapter was built to see: network congestion.
